\documentclass[10pt,journal,twoside]{IEEEtran} % for editor use only for final preprint

% For initial draft
%\documentclass[10pt,draftcls,journal,onecolumn]{IEEEtran}
%\usepackage{endfloat}
%\usepackage[switch]{lineno}
%\usepackage{lineno}


\usepackage{cite}
\usepackage{amsmath,amssymb,amsfonts}
\usepackage{graphicx}
\usepackage{siunitx}
\usepackage[colorlinks=true,allcolors=blue]{hyperref}
\usepackage{cleveref}
\crefname{equation}{}{}
\Crefname{equation}{}{}
\crefname{figure}{Fig.}{Figs.}
\Crefname{figure}{Fig.}{Figs.}
\crefname{table}{Table}{Tables}
\Crefname{table}{Table}{Tables}
\usepackage{booktabs}
\usepackage{multirow}

\usepackage{svg}
\svgpath{{./figures}}

\title{Detecting flaws in metal by mapping equipotentials on a two-dimensional graphite model} % Your title goes here
\author{Callie Butash\thanks{Author for correspondence: 426cbutash@frhsd.com}, Emily Chen, Jake Chin, Jason Donovan Katz, Grace Nealon, and Petra Rofman\thanks{Authors are with the Science \& Engineering Magnet Program, Manalapan High School, 20 Church Lane, Englishtown, NJ 07726, USA}} % Authors names here
% First author is normally a special position
% Corresponding author is a special position
% Last author normally most senior
% You can put your names in any order you decide but be aware of these conventions
\date{April 10, 2026}
\markboth{Journal of Science \& Engineering, Vol.~2, No.~4,~April 10, 2026}{Butash \MakeLowercase{\textit{et al.}}: Detecting flaws by mapping equipotentials} % may need shortened form of author names or title
\setcounter{page}{68} % LEAVE FOR EDITOR
\newcommand{\keywords}{equipotentials, two-dimensional, voltage, electric potential, electric field, graphite, three-dimensional (3-D), two-dimensional (2-D), non-destructive testing} % put your keywords here, lower case unless proper noun or acronym, separated by commas
\makeatletter
\AtBeginDocument{
\hypersetup{%
pdftitle={\@title},
pdfauthor={\@author},
pdfsubject={physics},
pdfkeywords={\keywords}}}
\IEEEoverridecommandlockouts
\IEEEpubid{\makebox[\columnwidth]{\href{https://10.64808/xxxxxx}{10.64804/xxxxxxxx}~\copyright~2026 J S\&E \href{https://creativecommons.org/licenses/by-nc/4.0/}{CC-BY-NC 4.0}\hfill} \hspace{\columnsep}\makebox[\columnwidth]{ }}% LEAVE FOR EDITOR
\makeatother







\begin{document}
%\linenumbers % for editor use
\maketitle

% AUTHORS PUT YOUR ABSTRACT HERE
\begin{abstract}
With the popularization of 3-D (three-dimensional) printed metals, the necessity for a method to detect internal flaws is increasing. We experimented with a 2-D (two-dimensional) graphite model of 3-D printed metal and measured the potential at various points throughout the model in a grid-like fashion. To simulate imperceptible flaws, we erased internal portions of the graphite and re-measured the potentials at the same points. Graphing our measurements, we created representations of nonlinear potential differences, allowing us to conclude that calculating potential in intervals over a conducting surface is a valid method to detect imperfections hidden within 3-D printed conducting surfaces under certain conditions.
\end{abstract}

\begin{IEEEkeywords}
\keywords
\end{IEEEkeywords}

\section{Introduction}
\IEEEPARstart{T}{he electric potential} of a particle can be described as the amount of electric potential energy per unit column of charge \cite{tipler,barrons,openstax}. This can be represented by the following equation:
\begin{equation}
\Delta V = - \int \vec{E} \cdot d\vec{r},
\label{eq:1}
\end{equation}
where $V$ is the potential difference that is represented by the integral of the electric field along a path, with the negative sign showing that electric potential will decrease in the direction of the field. Equipotential lines represent locations in an electric field where the potential is constant \cite{nave:2000:hyperphysics,knight:2017:physics}. We want to decipher how an electric field will differ when given a vertical or a horizontal gap utilizing measured potential and \cref{eq:1}. Previous research in this area has examined mapping equipotentials and electric fields for low-cost electrophoresis cells \cite{ayyagari:2024:computational,baru:2024:demonstrating,cohen:2024:visualizing,edwards:2024:electric,kumar:2024:mapping}, thus the technique has potential for our current application. 

As more capable additive manufacturing (i.e. 3D printing) technologies become more widely available, including in conductive materials such as metal, means of testing the integrity of parts nondestructively (non-destructive testing or NDT) are needed. Traditional industrial methods for NDT require specialized equipment (ultrasonic testing), radioactive sources (radiographic testing), or chemical dyes (penetrant testing), all of which would be unwieldy for small scale 3D printing processes \cite{griffiths:2013:introduction,kittel:2004:introduction}. Researchers have been finding defects in 3-D printed metal using x-ray imaging, heat cameras, and AI machine learning \cite{forrester:2023:researchers,loke:2013:recent,hauptmann:2018:revealing}. While their advanced technologies offer high-quality options to detect these flaws, our 2-D model is a more cost-effective and accessible option. Our solution utilizes potentials and the electric field and offers a quick quality check of the metals rather than an in-depth location search.

We hypothesize that blemishes in metal can be detected by observing a nonlinear voltage difference across a conducting surface. Since electric fields run perpendicular to equipotentials we expect that in the solid graphite region, the equipotentials will be linear, and perpendicular to the electric field. We also expect that at all grid points in the region with a vertical flaw will have nonlinear equipotential lines that run horizontally and curve around the discontinuity. Finally, we expect that the region with a horizontal flaw will have relatively linear equipotential lines. As the grid points get closer to where the flaw is located, the equipotential lines for this region may begin to curve around the discontinuity. 

To test these hypotheses, we set up an experiment to replicate, in two dimensions, the voltage across a conductive material: graphite. Our experiment models a block of metal as a graphite square on paper and flaws as erased portions. We conducted three trials: unflawed metal, metal with a vertical flaw, and metal with a horizontal flaw. 






\section{Methods and materials}
\subsection{Setup}
To test our hypotheses, we used a ruler to measure a box with dimensions \qtyproduct{50x50}{\milli\meter} on the corner of \qtyproduct{8.5x11}{inch} printer paper and marked along the length and width of the box in increments of \qty{10}{\milli\meter}. We then shaded the box as darkly and uniformly as possible with the graphite from a standard \#2 pencil (Dixon Ticonderoga; Appleton, WI). 

Two paper clips were straightened, placed along opposite sides of the box, and attached to a \qty{6}{\volt} power supply containing four \qty{1.5}{\volt} AA alkaline batteries (Energizer; Clayton, MO). The paperclips ensured good electrical contact. The negative lead from the battery pack was wrapped around the paperclip at the top of the square, and the positive lead was wrapped around the paperclip at the bottom. We used masking tape to secure the wires and paper clips to the paper surface.  

\begin{figure} 
\begin{center}
\includegraphics[width=\columnwidth]{figures/figure1.png}
\end{center}
\caption{Graphite square for control trial with grid}
\label{fig:1}
\end{figure}

\subsection{Data collection}
Using a digital multimeter (Fluke 106; Everett, WA), we measured the voltages with the negative probe on the negatively charged paper clip and the positive probe on the \qty{10}{\milli\meter} mark at the leftmost edge of the graphite square. Keeping the negative probe in place, we moved the positive probe to the right along the graphite in increments of \qty{10}{\milli\meter} until we reached the rightmost edge at \qty{50}{\milli\meter} and recorded the resulting voltage. We repeated this across each row of \qty{10}{\milli\meter} increments until we had voltage values for all intersections of \qty{10}{\milli\meter} grid lines, depicted as light blue dots in \cref{fig:1}.
  
\subsection{Horizontal and vertical flaws}

To mimic flaws in the metal, we erased a vertical line halfway through the box, spanning a length of \qty{45}{\milli\meter} and a width of approximately \qty{5}{\milli\meter}, as shown in \cref{fig:2}, to create a horizontal flaw. We repeated the aforementioned experimental process to record the voltage throughout the box in the same grid-like fashion. Then, we moved the paperclips to the two other sides of the box, causing the same erased ``flaw'' to be vertical, as shown in \cref{fig:3}, and repeated the process of measuring the voltage throughout the box in the same grid-like fashion.

\begin{figure}
\begin{center}
\includegraphics[width=\columnwidth]{figures/figure2.png}
\end{center}
\caption{Graphite square for trial with horizontal flaw}
\label{fig:2}
\end{figure}

\begin{figure}
\begin{center}
\includegraphics[width=\columnwidth]{figures/figure3.png}
\end{center}
\caption{Graphite square for trial with vertical flaw}
\label{fig:3}
\end{figure}

The flaw orientation is potentially important for NDT applications as different techniques are sensitive to flaws in particular orientations. For example, in traditional NDT, radiographic testing is most sensitive to flaws along the path of the x-rays; while ultrasonic testing is most sensitive to flaws the present a discontinuous interface normal to the beams \cite{griffiths:2013:introduction,kittel:2004:introduction}. 




\section{Results}
\subsection{Control}
In our first experiment with a solid region (setup from \cref{fig:1}), we had a fairly linear voltage difference across the graphite square. This is to be expected as the region is continuous, allowing the electric field to remain constant and the equipotential lines to be evenly spaced, as shown in \cref{fig:41,fig:42}.

The control trial demonstrates the relatively consistent potential measurements across each axis of the grid, providing a baseline measurement. If we didn’t know if the grid had flaws or not, this is what would be expected for a solid square. The equipotentials and field lines are all perpendicular to each other, with the constant spacing demonstrating a consistent electric field.
  
\begin{figure}
\begin{center}
\includesvg[width=\columnwidth]{newfig41.svg}
\end{center}
\caption{Potential difference across the uniform graphite grid}
\label{fig:41}
\end{figure}

\begin{figure}
\begin{center}
\includegraphics[width=\columnwidth]{figures/figure4-2.jpg}
\end{center}
\caption{Electric field (red) and equipotential lines (blue) across the uniform graphite grid}
\label{fig:42}
\end{figure}


\subsection{Horizontal flaw}
In the region with a horizontal flaw (\cref{fig:2}), we had a horizontal, linear voltage difference along the sides of the square with a center voltage of zero. This measurement lines up with the absence of graphite in that region. The graph of the potential difference (\cref{fig:51}) highlights the interruption of measured voltage in the region where the flaw is located. The diagram of equipotential and electric field lines (\cref{fig:52}) visualizes the effects of the gap on the grid. The distorted lines still fall perpendicular to each other, but in order to do so they bend around the void. 

\begin{figure}[h]
\begin{center}
\includesvg[width=\columnwidth]{newfig51.svg}
\end{center}
\caption{Potential difference across the horizontally flawed grid}
\label{fig:51}
\end{figure}

\begin{figure}[h]
\begin{center}
\includegraphics[width=\columnwidth]{figures/figure5-2.jpg}
\end{center}
\caption{Electric field (red) and equipotential lines (blue) across the horizontally flawed grid}
\label{fig:52}
\end{figure}

\subsection{Vertical flaw}
In the region with a vertical flaw (\cref{fig:3}), we had a vertical, linear voltage difference along the sides of the square with a center voltage of zero, again in the location where graphite was erased. The graph of the potential difference (\cref{fig:61}), however, fails to  highlight the interruption in voltage in the region where the flaw is located. The visualization of the voltage and electric field (\cref{fig:62}) demonstrates the distortion, which is noticeably less than in the horizontal flaw (\cref{fig:52}). 
  
\begin{figure}[h]
\begin{center}
\includesvg[width=\columnwidth]{newfig61.svg}
\end{center}
\caption{Potential difference across the vertically flawed grid}
\label{fig:61}
\end{figure}

\begin{figure}[h]
\begin{center}
\includegraphics[width=\columnwidth]{figures/figure6-2.jpg}
\end{center}
\caption{Electric field (red) and equipotential lines (blue) across vertically flawed grid}
\label{fig:62}
\end{figure}










\section{Discussion}
\subsection{Did the flaw cause a predictable effect on the measured potential values?}
In this experiment, we successfully measured and mapped potential values across flawed and unflawed graphite squares to mimic flawed and unflawed 3-D printed metals. The voltage measurements change smoothly throughout the graphite squares, except when the measurements near the gaps.

Comparing the graphs of the potential difference, however, shows that only the horizontally flawed trial produced a graph visibly different from the control trial. \Cref{fig:52,fig:62} show how the orientation of the flaw affects the potential differences over the region. The potential with the horizontally flawed trial drops significantly after the flaw, as the imperfection in the surface increases resistance and, in turn, leads to a reduced voltage. Ultimately, when the flaw is parallel to the equipotentials, and perpendicular to the electric field, as in the horizontal setup of \cref{fig:62}, current is much more disrupted, leading to more signal. 

Our findings confirm that flaw orientation is important for NDT applications, similar to how traditional ultrasonic testing is most sensitive to flaws the present a discontinuous interface normal to the beams \cite{griffiths:2013:introduction,kittel:2004:introduction}. 

\subsection{Can we generalize our findings to find flaws in 3D printed metal?}
        Similarly to our 2-D model, defects in 3-D printed metal will cause varying paths in the electric field. Our findings lead us to believe that electrical testing can be generalized to identify if a 3-D printed conductor has flaws that are not visible on the surface. While the vertical flaw was more difficult to detect than the horizontal one, there was still a gap in the potential measurements. In a similar manner to the 2-D experiment, a 3D object can have a voltage source connected across it and have its voltage measured at varying points. Analysis of the voltage will find flaws located where there are abrupt changes. However, these items have an entirely new dimension introduced where the electric field lines can stray, meaning that it may be more difficult to identify these imperfections using potential than in 2-D due to the increased number of data points that would need to be collected to form an adequate grid. Furthermore, we believe that the difficulties we faced identifying flaws in different orientations in the 2-D model will be amplified. 

\subsection{Sources of experimental error}
        Inconsistent application of graphite to the box could change the resistance, where a thicker graphite patch would have a lower resistance and a thinner graphite patch would have a higher resistance, ultimately resulting in the voltage drop not being perfectly linear. Additionally, if the line wasn’t perfectly erased, some current could still cross the remaining graphite, causing variations in the measured voltage. A last possible source of experimental error could be due to taking voltage measurements at inaccurate grid points; when eyeballing the grid points, the measurements may be slightly off their intended marks.

\section{Acknowledgment}
We thank the Science \& Engineering program and several anonymous reviewers. CB created figures for lab setup and results, wrote the Discussion section and developed the Methods and Materials before the experiment. EC created the line charts for the results and wrote portions of the Introduction, Results, and Sources of Experimental Error sections. JC wrote portions of the Introduction, abstract, and Methods and Materials section. JK followed instructions by assisting and gathering measurements for the first trial, taking over for the other two, and making results diagrams. GN set up the experiment and took measurements for the first trials. She also crafted the Results section and worked on the Introduction. PR wrote the Methods and Materials section, as well as part of the abstract, Results, and Discussion sections. All members were able to help others with writing, revising, and discussing data to formulate conclusions.




%Bibliography 
%[1] Georgia State University, Department of Physics and Astronomy (n.d.). Equipotential Lines. Hyper Physics. http://hyperphysics.phy-astr.gsu.edu/hbase/electric/equipot.html  
%[2] Forrester, Nikki. “Researchers Unveil New AI-Driven Method for Improving Additive Manufacturing | Argonne National Laboratory.” Argonne National Laboratory, 2023, www.anl.gov/article/researchers-unveil-new-aidriven-method-for-improving-additive-manufacturing. 
%\end{thebibliography}
\bibliographystyle{IEEEtran}
\bibliography{lab.bib}






% Images for the biography should be a 1in wide by 1.25 high or a 5:4 aspect ratio
% Keep your biography short and to the point or the editors will shorten it
\begin{IEEEbiography}[{\includegraphics[width=1in,height=1.25in,clip,keepaspectratio]{figures/butash.jpeg}}]{Callie Butash} is a senior in the Science and Engineering Magnet Program at Manalapan High School. She was an intern in the Monmouth County Engineering Department. When she is not studying for AP Physics C E\&M she enjoys sports and building a Spiderman-inspired grapple gun. 
\end{IEEEbiography}

\begin{IEEEbiography}[{\includegraphics[width=1in,height=1.25in,clip,keepaspectratio]{figures/chen.jpeg}}]{Emily Chen} is a senior in the Science and Engineering Magnet Program at Manalapan High School. She is a member of the Drama Club. She is currently an intern at ACON Pharmaceuticals in Cranbury, NJ. When she is not studying for AP Physics C E\&M she enjoys microspheres, musical theater, and working on her robotic fish. 
\end{IEEEbiography}

\begin{IEEEbiography}[{\includegraphics[width=1in,height=1.25in,clip,keepaspectratio]{figures/chin.jpeg}}]{Jake Chin} is a senior in the Science and Engineering Magnet Program at Manalapan High School. He is a member of the Drama Club. He is currently an intern with Girl in Space Club and is developing web-based applications to teach students astrophysics and orbital mechanics of satellite swarms. 
\end{IEEEbiography}

\begin{IEEEbiography}[{\includegraphics[width=1in,height=1.25in,clip,keepaspectratio]{figures/katz.jpeg}}]{Jason Donovan Katz} is a senior in the Science and Engineering Magnet Program at Manalapan High School. He is an avid tuba player and is president of the Beard Club. When he is not studying for AP Physics C E\&M he enjoys building robotic hands to play rock paper scissors. 
\end{IEEEbiography}

\begin{IEEEbiography}[{\includegraphics[width=1in,height=1.25in,clip,keepaspectratio]{figures/nealon.jpeg}}]{Grace Nealon} is a senior in the Science and Engineering Magnet Program at Manalapan High School. She was an intern in the Traffic Department at Monmouth County Engineering. She is the current reigning All Shore Irish Dancing Champion. In her free time, she enjoys singing musical numbers at S\&E Art and Science Night and building spy cars. 
\end{IEEEbiography}

\begin{IEEEbiography}[{\includegraphics[width=1in,height=1.25in,clip,keepaspectratio]{figures/rofman.jpeg}}]{Petra Rofman} is a senior in the Science and Engineering Magnet Program at Manalapan High School. She is a member of the Drama Club. She is currently an intern at ACON Pharmaceuticals in Cranbury, NJ. When she is not studying for AP Physics C E\&M she enjoys microspheres, origami, and working on her robotic fish. 
\end{IEEEbiography}
\vfill
\end{document}

